% ==================================================================
% id:       lim-016
% topic:    גבולות ורציפות של פונקציות
% subtopic: משפטי רציפות (ערך ביניים, ויירשטראס)
% solution: full
% ==================================================================
ביום ראשון בשעה
8:00
נכנס דניס למעלית בקומת הקרקע, ועלה בה עד לקומה
ה-$17$,
שאליה הגיע בדיוק בשעה
8:02.
למחרת, ביום שני,
נכנס דניס למעלית
בשעה
8:00
וירד בה עד קומת הקרקע, שאליה הגיע בדיוק בשעה
8:02.
נניח שבשני הימים הגובה של דניס מעל פני הקרקע משתנה באופן רציף עם הזמן, אך אין להניח שהמעלית נעה במהירות קבועה.

הוכיחו שקיים רגע בזמן  בין
8:00
ל-8:02
בו דניס חווה “דה־ז׳ה־וו”: כלומר, בשני הימים, בדיוק באותה שעה שמראה השעון, הוא נמצא באותו גובה בדיוק.
\hint{הגדירו פונקציית עזר: ההפרש בין הגובה ביום ראשון לגובה ביום שני באותה שעה, ובדקו את הסימן שלה בקצוות.}
\begin{solbox}
נסמן את הזמן בדקות שעברו מאז
8:00
ב-
$t\in[0,2]$.
נסמן ב-
$h_1(t)$
את גובהו של דניס ביום ראשון, וב-
$h_2(t)$
את גובהו ביום שני, באותו הזמן שמראה השעון. לפי הנתון, שתי הפונקציות רציפות בקטע
$[0,2]$.
כמו כן,
$$,h_1(0)=0,\quad h_1(2)=17$$
ואילו
$$.h_2(0)=17,\quad h_2(2)=0$$
נגדיר
$$.g(t)=h_1(t)-h_2(t)$$
אז
$g$
רציפה ב-
$[0,2]$,
ומתקיים
$$,g(0)=h_1(0)-h_2(0)=-17<0$$
וכן
$$.g(2)=h_1(2)-h_2(2)=17>0$$
לכן, ממשפט ערך הביניים, קיים
$c\in(0,2)$
כך ש-
$g(c)=0$.
כלומר
$$,h_1(c)=h_2(c)$$
וזה בדיוק אומר שבאותה שעה שמראה השעון, בשני הימים, דניס נמצא באותו גובה.
\end{solbox}
